\scanpage{51}
% archive sheet number: 50

\mspage{15}
\diagram{\begin{tikzcd}[ampersand replacement=\&]
\& \gamma(0,1)\arrow[dl,"j_0"'] \&\\
\gamma(0,0) \& \& \gamma(1,1)\arrow[ul,"i_{1}"']\arrow[dl,"j_1"]\\
\& \gamma(1,0)\arrow[ul,"i_0"'] \&
\end{tikzcd}}

\margin{De plus, on doit imposer autant pour les dérivateurs induits $\underline D_I$ $(I\in\mathrm{Diag})$, \unsure{idem|donc} dans $\underline D(I\times\Delta^1\times\Delta^1)$.}

\[
\text{alors}\quad
\left\{\begin{aligned}i_1\ \text{iso}&\Longleftrightarrow i_0\ \text{iso},\\
j_1\ \text{iso}&\Longleftrightarrow j_0\ \text{iso}.\end{aligned}\right.
\]

Cette dernière assertion résulte de la précédente (car cart. implique $\Rightarrow$, cocartésien implique $\Leftarrow$), pourvu qu'on sache ceci :

``\uline{Lemme}'' \deleted{(associé à Der 6)} : si $\xi=\deleted{\unsure{\xi}}$

\diagram{$\left(\begin{tikzcd}[ampersand replacement=\&,row sep=small]
\& \xi_a\arrow[dl,"j_0"']\\
\xi_0 \&\\
\& \xi_b\arrow[ul,"i_0"']
\end{tikzcd}\right)$}

est un diagramme substantiel tel que $i_0$ soit un iso, alors il est isom. à l'image inverse d'un $y\in\underline D(\underline U\deleted{\unsure{\Phi}})$ (où $\underline U=(\deleted{\unsure{b}}0\longrightarrow a)\simeq\Delta^1$ est l'ouvert canonique gauche de $\underline\Phi$) par la rétraction \inserted{$p$} de $\underline\Psi$ sur $\underline U$.

% margin note written sideways
\margin{\deleted{\unsure{cet axiome est lié à|cet axiome est le}}}

Il suffit de prouver que

\[
\deleted{\unsure{\xi}}\ p^*p_*(\xi)\longrightarrow\xi\quad\text{, adj.}
\]

est iso. Comme $p$ est bien sûr lisse, comme induit par $\mathrm{pr}_1\colon\Delta^1\times\Delta^1\longrightarrow\Delta^1$, sur un ouvert, mais aussi \deleted{propre} $p_*(\xi)$ se calcule fibre par fibre. On utilise que \deleted{\unsure{dans|$\leftarrow$}}

\[
A\simeq\underline D(e)\longrightarrow\underline D^{\unsure{\ell c|ec|cc}}(\Delta^1)
\]

est une équiv. de cat. (\uline{forme de Der 3})$^*$.
